\chapter{Conclusiones y trabajo futuro}
\label{chap:conclusiones}

\ph{El último capítulo cierra la memoria. Se divide habitualmente en dos
secciones: las \emph{conclusiones} (qué se ha conseguido) y el \emph{trabajo
futuro} (qué queda fuera y se podría abordar a continuación).}

\section{Conclusiones}
\label{sec:conclusiones}

\ph{Estructura recomendada:}

\ph{1. \textbf{Recapitulación del problema}: una o dos frases recordando
qué problema se planteaba al inicio (sin entrar en detalles, ya que el
lector viene del capítulo de resultados).}

\ph{2. \textbf{Cumplimiento de los objetivos}: revisa los objetivos
específicos planteados en el Capítulo~\ref{chap:intro} y argumenta cómo
se han cubierto. Una buena estrategia es dedicar un párrafo a cada
objetivo, en el mismo orden en que se enunciaron.}

\begin{itemize}
  \item \ph{El \textbf{primer objetivo} \dots se ha cumplido mediante \dots.}
  \item \ph{El \textbf{segundo objetivo} \dots se ha cumplido mediante \dots.}
  \item \ph{El \textbf{tercer objetivo} \dots se ha cumplido mediante \dots.}
  \item \ph{El \textbf{cuarto objetivo} \dots se ha cumplido mediante \dots.}
  \item \ph{El \textbf{quinto objetivo} \dots se ha cumplido mediante \dots.}
\end{itemize}

\ph{3. \textbf{Aportaciones más amplias}: comenta brevemente la contribución
del trabajo más allá del producto en sí. Por ejemplo, aplicación de un
paradigma a un dominio nuevo, integración con un ecosistema, validación
empírica de una técnica, etc.}

\ph{4. \textbf{Reflexión sobre la planificación}: comenta la desviación
final de horas, en qué fases se concentró y qué aprendizajes se extraen.}

\ph{5. \textbf{Cierre}: una frase que cierre el capítulo conectando las
limitaciones identificadas en la introducción con la solución desarrollada
(``el sistema desarrollado resuelve los X problemas estructurales\dots'').}

\section{Trabajo futuro}
\label{sec:trabajo_futuro}

\ph{Líneas de extensión del trabajo, ordenadas de mayor a menor inmediatez.
Cada línea se redacta como un párrafo corto que explique \emph{qué} se
podría hacer y \emph{por qué} sería interesante.}

\textbf{\ph{Línea 1: título corto.}} \ph{Descripción: qué se podría añadir,
por qué encaja con el sistema actual y qué dificultad supondría.}

\textbf{\ph{Línea 2: título corto.}} \ph{Descripción.}

\textbf{\ph{Línea 3: título corto.}} \ph{Descripción.}

\textbf{\ph{Línea 4: título corto.}} \ph{Descripción.}

\textbf{\ph{Línea 5: título corto.}} \ph{Descripción.}
